\documentclass[a4]{article}
%\textwidth=4.5in
\textwidth=14cm
%\textheight=7.125in
\textheight=22cm
%\usepackage{pseudocode} %sak
\usepackage{funpseudocode}
\usepackage[pdfpagemode={UseOutlines},bookmarks=true,bookmarksopen=true,
   pdfauthor={S. Arun-Kumar},pdfcreator={S. Arun-Kumar},
   bookmarksopenlevel=0,bookmarksnumbered=true,hypertexnames=false,
   colorlinks,linkcolor={blue},citecolor={red},urlcolor={magenta},
   pdfstartview={FitV},unicode,breaklinks=true]{hyperref}
   


\hyphenation{pa-ram-e-ters sep-a-rated}
%
%
%    September 2, 1999
%
%
\title{FunPseudocode: A \LaTeX\  Style File for Displaying \underline{All} Algorithms}
\author{
D.L.\ Kreher\\
Department of Mathematical Sciences\\
Michigan Technological University \\ 
Houghton, MI 49931\\
{\tt kreher@mtu.edu}\\[1ex]
and \\[1ex]
D.R.\ Stinson\\
Department of Combinatorics and Optimization\\
University of Waterloo\\
Waterloo ON, N2L 3G1\\
{\tt dstinson@uwaterloo.ca}\\[1ex]
and (last and least)\\[1ex]
S. Arun-Kumar\\
Department of Computer Science and Engineering\\
Indian Institute of Technology Delhi\\
Hauz Khas, New Delhi 110016, India.\\
{\tt sak@cse.iitd.ac.in}
}

\newcommand{\keyword}[1]{\texttt{#1}}
\begin{document}
%\makeatletter
%\setlength{\pcode@width}{2.75in}
%\addtolength{\pcode@width}{-40pt}
%\makeatother

\maketitle
\begin{abstract}
  There are many styles of programming and correspondingly there 
  ought to be just as many styles of specifying algorithms in pseudocode.
  Due to various historical reasons there is a ``von Neumann 
  stranglehold'' (deriving its power mainly from the ``von
  Neumman bottlneck''\cite{backus}) which dominates the domain of pseudocode. 

  While any number of styles and style files exist to format algorithms written
  in \LaTeX\ in the imperative style, there are none for writing algorithms in
  `functional pseudocode'. The last author has decided to make amends in this
  document -- with minimal modification of the {\tt pseudocode.sty} file, mainly
  because this was the simplest and easiest thing to do. Most of the other style
  files looked more complicated and harder to understand and tweak. The new
  style file is called \keyword{funpseudocode.sty} and its commands are used in
  the environment \keyword{pseudocode}. 
\end{abstract}

\section{Introduction}

This paper describes a \LaTeX\ environment named \keyword{pseudocode}
that can be used for describing algorithms in pseudocode form.  This
is the style used in the %our %recently published 
textbook {\em Combinatorial Algorithms: Generation, Enumeration and Search}
  \cite{cages} written by the first two authors.
The style file \keyword{pseudocode.sty} is freely available on

\url{http://www.math.mtu.edu/~kreher/cages.html}
while the style file \keyword{funpseudocode.sty} is available from the last
author on request.


This package is quite easy to use, and allows imperative (i.e. state-based)
algorithms to be described in a \LaTeX\ document using a natural Pascal-like
syntax and functional algorithms (which are not state-based) in a simple
mathematical notation that is more precise (and less ambiguous) than in most
mathematics texts.  In the remaining sections of this note, we describe how to
use the \keyword{pseudocode} environment and we present some examples.  Readers
familiar with \LaTeX\ (see \cite{lamport}) should be able to easily customize
the style file to include additional desired features.

The \keyword{pseudocode} environment requires the \keyword{fancybox}
package by Timothy Van Zandt.  This package is described in 
Section 10.1.3 of \cite{LC}.
Other environments for describing algorithms include
\keyword{alg}, \keyword{algorithmic}, \keyword{newalg} and \keyword{program}.
These style files, as well as \keyword{fancybox}, are all available from the CTAN web site
\verb+http://www.ctan.org/+

In this paper we describe our extension to the \texttt{pseudocode}
package to include pseudocode written in the functional and the
declarative styles as well.

\paragraph{Imperative pseudocode} The package \texttt{pseudocode.sty}
written by the first two authors aids the presentation of essentially
imperative algorithms in pseudocode that could be easily expanded or
refined into actual code written in a high-level block-structured
language. We call this {\em imperative pseudocode}.

There are other styles of programming which are not necessarily state-based and which require algorithms to be
specified in pseudocode. Functional and declarative algorithms
do not necessarily suffer from the complexities of sequencing
low-level state updates. Indeed much of pure functional programming
does not require the notion of state. More importantly, state-based
algorithm specification may actually hinder the development of code in
a functional or declarative programming language and make it harder to
understand and code even fairly easy problems. We call this the ``von
Neumann stranglehold''. It actually prevents undergraduate students
trained to program in imperative languages from
\begin{itemize}
\item relating algorithm design and programming to their mathematical
  education and
\item prevents clarity of thought and precision in expression when it
  comes to specifying recursion in data and programs. It has led to a
  colleague commenting that the ``for-loop is the most basic instinct 
  after sex''.
\end{itemize}

Functional specifications are often more succinct than state-based
specification.  The same also holds for algorithms written in a
declarative style. Logically speaking there is no reason why an
algorithm cannot be specified as a set of rules in a first- or
higher-order logical language.

\paragraph{Functional pseudocode} As the name suggests, it is simply a
way of specifying and writing algorithms as mathematical
functions. While most recursive functions in any recursive function
theory text are actually algorithms, one may ask why it is necessary to have a
pseudocode meta-language to describe functional algorithms. There are many
reasons for this.
\begin{itemize}
\item Algorithms written as mathematical functions using the usual devices
  available for formatting mathematics, pay scant regard to information about
  the scope of names and the life-times of variables and functions. Functions
  may be deeply nested or localised in other functions and expressions and
  normal mathematical formatting would ignore the nesting of definitions. We
  will see examples of this as we go along.
\item Mathematically expressed algorithms do not allow for the specification of
  orders of evaluation which are often of interest in the case of
  algorithms. For instance, the definition of a function by cases such as
  $|x| = \left\{ 
    \begin{array}{ll} 
      x & \mbox{~if~} x \geq 0\\
      -x & \mbox{~if~} x < 0
    \end{array}
    \right.$
  does not take into account that the condition would be first evaluated. It
  also does not take into account that the condition needs to be  evaluated only
  once for both cases. A functional pseudo-code definition as given below
 
  \begin{pseudocode}{abs}{x}
  \BEGIN
    \CASES
    \IT{x \geq 0}{x}
    \IT{\ELSE}{-x}
    \ENDCASES
  \END
\end{pseudocode}

allows these distinctions to be taken into account.
\item Functional pseudo-code permits an almost direct translation of the
  algorithm into a program (in a linearised form) in a modern functional
  programming language.
\end{itemize}

Such issues are important because they greatly affect our proofs of correctness
and our calculations of the efficiency of algorithms expressed as
functions. Further mathematicians, logicians and programmers are in general
sloppy about distinguishing between various forms of equality and
equivalence\footnote{C, C++, Java and python programmers are even worse! For
  instance, {\tt x = x+1} must easily rank as the most bizarre equation ever
  encountered by an unsuspecting pure mathematician. It collapses the entire set
  of naturals to a singleton!}.  In particular, it is important to distinguish
between an equality by definition and an equality derived by other means.


\paragraph{Declarative pseudocode} Arguments similar to those in favour of
functional pseudocode may also be advanced in favour of expressing
algorithms as merely logical rules. While in a functional algorithm
there are clear and distinct notions of input to a function and its 
output, the declarative style is more relational in its approach and
therefore it should suffice to specify a collection of logical rules 
which make up an algorithm. 

\paragraph{Mixed pseudocode} There are however many algorithms that
are expressed in a mixture of styles, just as no programming language
is either purely functional, or purely declarative or purely
imperative. It is our contention that for each problem or sub-problem,
any or a mixture of styles of expressing algorithms may be used,
depending upon taste, ``natural-ness'', conciseness, efficiency etc.

\paragraph{Annotations} A final stylistic improvement (independent of the
programming paradigm) is that in this environment we allow comments to specify
preconditions, post-conditions, invariants, bound functions etc to simplify the
understanding of the correctness of the algorithm described in the pseudo-code

Since this package unifies the existing imperative style with the functional and
declarative styles, it has been found necessary to explain all aspects including
that of the imperative style. What better way to do this than to extend the
original documentation written by the first two authors for this purpose? That
is precisely what we do in this document.


Without any further ado we plunge into the matter of the original document and
intersperse it with our own material. We hope the reader gets a holistic
perspective of algorithm specification and ``Houdini'' out of the von Neumann
stranglehold.

\section{The \keyword{pseudocode} Environment}

Within the \keyword{pseudocode} environment, a number of commands for 
popular algorithmic constructs are available.
In general, the commands provided can be nested to
describe quite complex algorithms. 

The \keyword{funpseudocode} environment is  invoked
as follows:
\begin{verbatim}
\begin{pseudocode}{<Name>}{<Parameters>}
  pseudocode constructs
\end{pseudocode}
\end{verbatim}

\noindent{}The argument \verb+<Name>+ is the name of the algorithm,
and \verb+<Parameters>+ is a list of parameters for the algorithm.
For example, the commands

%\iffalse{ %sak
\begin{verbatim}
\begin{pseudocode}{CelsiusToFahrenheit}{c}
  f \GETS {9c/5} + 32\\
  \RETURN{f}
\end{pseudocode}
\end{verbatim}
%}\fi %sak

\noindent{}produces the following output when included in
a \LaTeX\  document:

\begin{pseudocode}{CelsiusToFahrenheit}{c}
  f \GETS {9c/5} + 32\\
  \RETURN{f}
\end{pseudocode}

%\iffalse{ %sak
Notice that the command \verb+\GETS+ produces the Pascal assignment
symbol, which we use to indicate an assignment of a
variable\footnote{The first two authors in their original package used
  a left arrow for assingment. The last author has changed it to the
  Pascal-like assignment symbol since the left arrow is more useful
  specifying Horn clauses in a declarative algorithm}. The user
could use instead some other symbol, if desired.  For example,
\verb+\GETS+ could be replaced by \verb+=+, as is done in the ``C''
programming language.
%}\fi %sak 

The corresponding functional algorithm would read in  \LaTeX\ as
follows.

\begin{verbatim}
\begin{pseudocode}{CelsiusToFahrenheit}{c}
  {9c/5} + 32\\
\end{pseudocode}
\end{verbatim}

\noindent{}and produce the following output.

\begin{pseudocode}{CelsiusToFahrenheit}{c}
%%  f \GETS {9c/5} + 32\\ %sak
  {9c/5} + 32
%%  \RETURN{f}
\end{pseudocode}

Here notice that we use ``$\DEF$'' to denote ``equality by
definition''.  Every function definition will use this symbol to
separate the name and formal parameters of the function from its
body. We use the standard equality symbol ``$=$'' to denote the
usual notion of equality used in (equational) identities.

\subsection{ The {\em begin-end} Construct}

To form compound statements from simple statements, the {\em begin-end} 
construct is used as follows:
\begin{verbatim}
 \BEGIN
    some statement\\
    another statement\\
    yet another statement
 \END 
\end{verbatim}
This generates the following:

\medskip
\begin{pseudocode}[display]{}{}
\BEGIN
\mbox{some statement}\\
\mbox{another statement}\\
\mbox{yet another statement}
\END
\end{pseudocode}
\\[2ex]
\noindent{}The effect of this construct is to group 
a collection of statements
using a left brace bracket of the appropriate size.

In the sections that follow we will use the notation
\verb+<stmt>+
to indicate  a simple statement or a compound statement.
Note that the contents of statements are typeset in math mode.
Therefore, non-math mode text must be enclosed in an
\verb+\mbox{}+.

Observe that the double backslash \verb+\\+ plays the same
role as the semicolon in Pascal, i.e., it is used to
separate statements, and should never appear before
\verb+\END+.

%\iffalse{ %sak
\subsection{The {\em if-then-else} Construct}

The {\em if-then-else} construct takes various forms,
such as the following:
\begin{verbatim}
\IF <condition> \THEN <stmt> 
\IF <condition> \THEN <stmt> \ELSE <stmt> 
\IF <condition> \THEN <stmt> \ELSEIF <stmt> \THEN <stmt> 
\end{verbatim}
Note that there is no limit placed on the number
of \verb+\ELSEIF+s that may be used in an {\em if-then-else} construct.
For example,  the commands:
\begin{verbatim}
\IF some condition is true
\THEN
 \BEGIN
    some statement\\
    another statement\\
    yet another statement
 \END
\ELSEIF some other condition is true
\THEN
 \BEGIN
    some statement\\
    another statement\\
    yet another statement
 \END
\ELSEIF some even more bizarre condition is met
\THEN 
  do something else
\ELSE
  do the default actions
\end{verbatim}

\noindent{}would produce the following output:

\medskip
\begin{pseudocode}[display]{}{}
\IF \mbox{some condition is true}
\THEN 
\BEGIN
\mbox{some statement}\\
\mbox{another statement}\\
\mbox{yet another statement}
\END
\ELSEIF \mbox{some other condition is true}
\THEN 
\BEGIN
\mbox{some statement}\\
\mbox{another statement}\\
\mbox{yet another statement}
\END
\ELSEIF\mbox{some even more bizarre condition is met}
\THEN \mbox{do something else}
\ELSE
\mbox{do the default actions}
\end{pseudocode}

\subsection{The {\em for} Loop}

The {\em for} loop takes the following forms:
\begin{verbatim}
\FOR <var> \GETS <lower>  \TO <upper> \DO <stmt> 
\FOR <var> \GETS <upper>  \DOWNTO <lower> \DO <stmt> 
\FOREACH <condition> \DO  <stmt> 
\end{verbatim}
For example,
\begin{verbatim}
\FOR i \GETS 0 \TO 10 \DO
  some processing
\end{verbatim}
produces

\medskip
\begin{pseudocode}[display]{}{}
\FOR i \GETS 0 \TO 10 \DO
\mbox{some processing}
\end{pseudocode}

\noindent{}and

\begin{verbatim}
\FOREACH x \in \mathcal{S}  \DO
  some processing 
\end{verbatim}

\noindent{produces}

\medskip
\begin{pseudocode}[display]{}{}
\FOREACH x \in \mathcal{S}  \DO
  \mbox{some processing}
\end{pseudocode}

\subsection{The {\em while} Loop}

The {\em while} loop takes the following form:

\begin{verbatim}
\WHILE <condition> \DO <stmt> 
\end{verbatim}

\noindent{}For example,

\begin{verbatim}
\WHILE some condition holds \DO 
  some processing 
\end{verbatim}

\noindent{}produces

\medskip
\begin{pseudocode}[display]{}{}
\WHILE \mbox{some condition holds} \DO 
\mbox{some processing }
\end{pseudocode}

\subsection{The {\em repeat-until} Loop}

The {\em repeat-until} loop takes the following form:

\begin{verbatim}
\REPEAT <stmt> \UNTIL <condition> 
\end{verbatim}

\noindent{}For example,

\begin{verbatim}
\REPEAT
  some processing 
\UNTIL some condition is met 
\end{verbatim}

\noindent{}produces

\medskip
\begin{pseudocode}[display]{}{}
\REPEAT
\mbox{some processing}
\UNTIL \mbox{some condition is met}
\end{pseudocode}

\subsection{Main Programs and Procedures}

We can describe a main program that calls one (or more) procedures as follows:

\begin{verbatim}
\begin{pseudocode}{<Name>}{<Parameters>}
\PROCEDURE{<ProcedureName>}{<ProcedureParameters>}
   some stuff
\ENDPROCEDURE
\MAIN
  some stuff\\
  \CALL{<ProcedureName>}{<ActualParameters}>\\
  more stuff
\ENDMAIN
\end{pseudocode}
\end{verbatim}

Here is a simple example to illustrate the 
use of a main program calling a procedure.
The commands

\begin{verbatim}
\begin{pseudocode}{TemperatureTable}{lower, upper}
\PROCEDURE{CelsiusToFahrenheit}{c}
  f \GETS {9c/5} + 32\\
  \RETURN{f}
\ENDPROCEDURE 
\MAIN
x \GETS lower \\
\WHILE x \leq upper \DO
\BEGIN
 \OUTPUT{x, \CALL{CelsiusToFahrenheit}{x}}\\
 x \GETS x+1
\END
\ENDMAIN
\end{pseudocode}
\end{verbatim}

\noindent{}produce the following output:

\begin{pseudocode}{TemperatureTable}{lower, upper}
\PROCEDURE{CelsiusToFahrenheit}{c}
  f \GETS {9c/5} + 32\\
  \RETURN{f}
\ENDPROCEDURE 

\MAIN
x \GETS lower \\
\WHILE x \leq upper \DO
\BEGIN
 \OUTPUT{x, \CALL{CelsiusToFahrenheit}{x}}\\
 x \GETS x+1
\END
\ENDMAIN
\end{pseudocode}
%}\fi %sak

\subsection{Comments}

A comment statement may be inserted in an algorithm using the
following command:
\begin{verbatim}
\COMMENT{<commenttype>}{<stmt>}
\end{verbatim}
For example, the commands
\begin{verbatim}
expression1 
\COMMENT{comment}{Now evaluate expression1}\\
expression2 
\COMMENT{comment}{evaluate expression2}\\
\end{verbatim}

\noindent{}produce the output

\medskip
\begin{pseudocode}[display]{}{}
expression1  
\COMMENT{comment}{Now evaluate expression1}\\
expression2  
\COMMENT{comment}{evaluate expression2}\\
\end{pseudocode}

\medskip

Note that comments are assumed to be text. Thus, in order to include
mathematical expressions in a comment, math mode must be used explicitly. Further, note that there are {\em two} parameters in a comment. This gives the designer freedom to classify comment types of various kinds. One common use is to allow the specification of preconditions, post-conditions, loop invariants and bound functions before loops and for tail-recursive functions. The following imperative algorithm illustrates these uses.

\begin{pseudocode}{Factorial}{n}
  \BEGIN
     \COMMENT{requires}{integer $n > 0$}\\
     \LOCAL{p \GETS 1}\\
     \FOR i  \GETS 1 \TO n \DO\\
     \BEGIN 
        p \GETS p*i\\
        \COMMENT{inv}{$0 < i \leq n \wedge p = i!$}\\
        \COMMENT{bf}{$n-i$}\\ 
     \END\\
     \COMMENT{ensures}{$p = n!$}\\
     \RETURN{p}
  \END
\end{pseudocode}




\subsection{A Declarative Example}
Consider the following logic program (as opposed to Prolog) lifted
straight from Llyod's book \cite{lloyd}. A logic program is simply an
unordered collection of Horn clauses. However the sequence below shows
a top-down development of the logic program.

\begin{pseudocode}{sort}{x,y}
  \CASES
  \HORN{sort(x,y)}{perm(x,y),sorted(x,y)}
  \HORN{sorted(nil)}{}
  \HORN{sorted(x.nil)}{}
  \HORN{sorted(x.y.z)}{lesseq(x,y),sorted(y,z)}
  \HORN{lesseq(x,x)}{}
  \HORN{lesseq(x,y)}{x < y}
  \HORN{perm(nil, nil)}{}
  \HORN{perm(x.y, u.v)}{delete(u, x.y, z), perm(z, v)}
  \HORN{delete(x, x.y, y)}{}
  \HORN{delete(x, y.z, y.w)}{delete(x, z, w)}
  \ENDCASES
\end{pseudocode}

\noindent{}generated from the following \LaTeX\ source 

\begin{verbatim}
\begin{pseudocode}{sort}{x,y}
  \HORN{sort(x,y)}{perm(x,y),sorted(x,y)}
  \HORN{sorted(nil)}{}
  \HORN{sorted(x.nil)}{}
  \HORN{sorted(x.y.z)}{lesseq(x,y),sorted(y,z)}
  \HORN{lesseq(x,x)}{}
  \HORN{lesseq(x,y)}{x < y}
  \HORN{perm(nil, nil)}{}
  \HORN{perm(x.y, u.v)}{delete(u, x.y, z), perm(z, v)}
  \HORN{delete(x, x.y, y)}{}
  \HORN{delete(x, y.z, y.w)}{delete(x, z, w)}
\end{pseudocode}
\end{verbatim}

\subsection{Some Functional Examples}

The following functional pseudo-code specifies the function for computing the
greatest non-negative integer no greater than the positive square root of an
integer. First, we give the specification in as a pure mathematical definition
(by cases if necessary) and then as functional pseudocode.  

\[
\lfloor \sqrt {n} \rfloor = 
\left\{
  \begin{array}{ll}
    undefined &\mbox{\ if\ } n < 0 \\
    k & \mbox{\ if\ } n \geq 0
  \end{array}
\right.
\]
where $0\leq k^2 \leq n < (k+1)^2$.

Since functions could be partial, we use ``$\bot$'' to stand for the undefined.
In any definition by
cases, it makes more sense to provide the condition to be satisfied before the
value, since it general the condition will be evaluated first. Further, the
``where'' clause is incorporated as part of the algorithm in the pseudocode.

\begin{pseudocode}{isqrt0}{n}
  \BEGIN
    \CASES
    \IT{n < 0}{\bot}
    \IT{n \geq 0}{k}
    \ENDCASES\\
  \WHERE 0 \leq k^2 \leq n < (k+1)^2      
  \END
\end{pseudocode}

\noindent{}whose  \LaTeX\ source is as follows.

\begin{verbatim}
\begin{pseudocode}{isqrt0}{n}
  \BEGIN
    \CASES
    \IT{\bot}{n < 0}
    \IT{k}{n \geq 0}
    \ENDCASES\\
  \WHERE 0 \leq k^2 \leq n < (k+1)^2      
  \END
\end{pseudocode}  
\end{verbatim}

In any definition by cases when the last case is implicitly a negation of the
cases preceding it, it may be more appropriate to use the keyword {\bf else}.


\begin{pseudocode}{isqrt0}{n}
  \BEGIN
    \CASES
    \IT{n < 0}{\bot}
    \IT{\ELSE}{k}
    \ENDCASES\\
  \WHERE 0 \leq k^2 \leq n < (k+1)^2      
  \END
\end{pseudocode}


Since most algorithms in a functional language would use several
conditional clauses which are reminiscent of Horn clauses in a logic
program, we have defined the keyword ``{\tt \textbackslash IT}''
(standing for ``if'' $\ldots$ ``then'' $\ldots$)
with two parameters separated by a ``$\rightarrow$''. The keywords
``{\tt \textbackslash CASES}'' and ``{\tt \textbackslash ENDCASES}''
enclose the various cases in a definition and align the
``$\rightarrow$'' symbols appropriately.

The following is an implementation obtained using top-down
refinement.  Notice the nesting of functions. Also notice that
\REC is a keyword in functional algorithms. \REC  indicates the presence of recursion in a function.

\begin{pseudocode}{isqrt1}{n}
  \BEGIN
    \BEGIN
    \CASES
    \IT{n < 0}{\bot}
    \IT{n = 0}{0}
    \IT{\ELSE}{\CALL{shrink}{n,0,n}}
    \ENDCASES
    \END\\
    \WHERE\\
    \REC \FUN {shrink} {n, l, u}\\
    \BEGIN
      \CASES
      \IT{l=u}{l}
      \IT{l<u \wedge (l+1)^{2} \leq n}{\CALL{shrink}{n, l+1, u}}
      \IT{l<u \wedge (l+1)^{2} > n}{l}
      \IT{l < u \wedge u^2 > n}{\CALL{shrink}{n, l, u-1}}
      \IT{l < u \wedge (u-1)^2 > n}{u-1}
      \IT{\ELSE}{\bot}\\
      \ENDCASES
     \END
     \ENDFUN
  \END\\
\end{pseudocode}

The above code was produced by the following \LaTeX\ code.

\begin{verbatim}
\begin{pseudocode}{isqrt1}{n}
  \BEGIN
    \BEGIN
    \CASES
    \IT{n < 0}{\bot}
    \IT{n = 0}{0}
    \IT{\ELSE}{\CALL{shrink}{n,0,n}}
    \ENDCASES
    \END\\
    \WHERE\\
    \REC \FUN {shrink} {n, l, u}\\
    \BEGIN
      \CASES
      \IT{l=u}{l}
      \IT{l<u \wedge (l+1)^{2} \leq n}{\CALL{shrink}{n, l+1, u}}
      \IT{l<u \wedge (l+1)^{2} > n}{l}
      \IT{l < u \wedge u^2 > n}{\CALL{shrink}{n, l, u-1}}
      \IT{l < u \wedge (u-1)^2 > n}{u-1}
      \IT{\ELSE}{\bot}\\
      \ENDCASES
     \END
     \ENDFUN
  \END\\
 \end{pseudocode}  
\end{verbatim}

We could equivalently have written this functional algorithm using a
\LET as follows.

\begin{pseudocode}{isqrt1}{n}
  \LET 
  \REC \FUN
  {shrink} {n, l, u}
   \BEGIN
     \CASES
      \IT{l=u}{l}
      \IT{l < u \wedge (l+1)^{2} \leq n}{\CALL{shrink}{n, l+1, u}}
      \IT{l < u \wedge (l+1)^{2} > n}{l}
      \IT{l < u \wedge u^2 > n}{\CALL{shrink}{n, l, u-1}}
      \IT{l < u \wedge (u-1)^2 > n}{u-1}
      \IT{\ELSE}{\bot}
      \ENDCASES
    \END\\
  \ENDFUN
   \IN
  \BEGIN
    \CASES
    \IT{n < 0}{\bot}
    \IT{n = 0}{0}
    \IT{\ELSE}{\CALL{shrink}{n,0,n}}
    \ENDCASES
  \END
 \end{pseudocode}

\noindent{}which has been produced by the following \LaTeX\ code.

\begin{verbatim}
\begin{pseudocode}{isqrt1}{n}
  \LET 
  \REC \FUN
  {shrink} {n, l, u}
   \BEGIN
     \CASES
      \IT{l=u}{l}
      \IT{l < u \wedge (l+1)^{2} \leq n}{\CALL{shrink}{n, l+1, u}}
      \IT{l < u \wedge (l+1)^{2} > n}{l}
      \IT{l < u \wedge u^2 > n}{\CALL{shrink}{n, l, u-1}}
      \IT{l < u \wedge (u-1)^2 > n}{u-1}
      \IT{\ELSE}{\bot}
      \ENDCASES
    \END\\
  \ENDFUN
   \IN
  \BEGIN
    \CASES
    \IT{n < 0}{\bot}
    \IT{n = 0}{0}
    \IT{\ELSE}{\CALL{shrink}{n,0,n}}
    \ENDCASES
  \END
 \end{pseudocode}  
\end{verbatim}

Rather than evaluate the condition $l < u$ repeatedly we could use
nested constructs such as the following.

\begin{pseudocode}{isqrt1}{n}
  \LET 
  \REC \FUN
  {shrink} {n, l, u}
   \BEGIN
     \CASES
      \IT{l=u}{l}
      \IT{l < u}{\BEGIN
                  \CASES
                  \IT{(l+1)^{2} \leq n}{\CALL{shrink}{n, l+1, u}}
                  \IT{(l+1)^{2} > n}{l}
                  \IT{u^2 > n}{\CALL{shrink}{n, l, u-1}}
                  \IT{(u-1)^2 > n}{u-1}
                 \ENDCASES
                 \END}
      \IT{\ELSE}{\bot}
      \ENDCASES
    \END\\
  \ENDFUN
   \IN
  \BEGIN
    \CASES
    \IT{n < 0}{\bot}
    \IT{n = 0}{0}
    \IT{\ELSE}{\CALL{shrink}{n,0,n}}
    \ENDCASES
  \END
 \end{pseudocode}  

The above pseudocode was produced by the following \LaTeX\ code.

\begin{verbatim}
\begin{pseudocode}{isqrt1}{n}
  \LET 
  \REC \FUN
  {shrink} {n, l, u}
   \BEGIN
     \CASES
      \IT{l=u}{l}
      \IT{l < u}{\BEGIN
                  \CASES
                  \IT{(l+1)^{2} \leq n}{\CALL{shrink}{n, l+1, u}}
                  \IT{(l+1)^{2} > n}{l}
                  \IT{u^2 > n}{\CALL{shrink}{n, l, u-1}}
                  \IT{(u-1)^2 > n}{u-1}
                 \ENDCASES
                 \END}
      \IT{\ELSE}{\bot}
      \ENDCASES
    \END\\
  \ENDFUN
   \IN
  \BEGIN
    \CASES
    \IT{n < 0}{\bot}
    \IT{n = 0}{0}
    \IT{\ELSE}{\CALL{shrink}{n,0,n}}
    \ENDCASES
  \END
 \end{pseudocode}  
\end{verbatim}

We could also pepper the code with preconditions, postconditions, invariants and
bound functions. In the above algorithm {\sc shrink} is a tail-recursive
function, that shrinks the interval (containing the integer-square-root) rather
slowly. A faster algorithm would be the function {\sc shrink2} shown in the
following pseudo-code. Since this is also less obvious algorithm, whose
correctness may be in doubt, it pays to include preconditions, post-constions and
invariants and give an annotated piece of pseudocode.

\begin{pseudocode}{isqrt2}{n}
  \COMMENT{precondition}{integer $n > 1$}\\
  \BEGIN
  \LET 
  \REC \FUN
  {shrink2} {n, l, u}
   \COMMENT{requires}{integers $n > 0 \wedge 0\leq l <  u\leq n$}\\
   \COMMENT{inv}{$0\leq l^2\leq n < u^2$}\\
   \COMMENT{bf}{$u-l-1$}\\
   \BEGIN
     \CASES
     \IT{u=l+1}{l}
     \IT{\ELSE}{\BEGIN
                   \CASES
                    \IT{m^2 > n}{\CALL{shrink2}{n, l, m}}
                    \IT{\ELSE}{\CALL{shrink2} {n, m, u}}
                    \ENDCASES
                   \END}
                 \WHERE m = (l+u)/2
      \ENDCASES
    \END\\
    \COMMENT{ensures}{$0\leq l^2\leq n<(l+1)^2$}
  \ENDFUN
   \IN
  \BEGIN
    \CASES
    \IT{n < 0}{\bot}
    \IT{n \leq 1}{n}
    \IT{\ELSE}{\CALL{shrink2}{n,0,n}}
    \ENDCASES
  \END
  \END
 \end{pseudocode}  

The above was produced by the following \LaTeX\ code.

\begin{verbatim}
\begin{pseudocode}{isqrt2}{n}
  \COMMENT{precondition}{integer $n > 1$}\\
  \BEGIN
  \LET 
  \REC \FUN
  {shrink2} {n, l, u}
   \COMMENT{requires}{integers $n > 0 \wedge 0\leq l <  u\leq n$}\\
   \COMMENT{inv}{$0\leq l^2\leq n < u^2$}\\
   \COMMENT{bf}{$u-l-1$}\\
   \BEGIN
     \CASES
     \IT{u=l+1}{l}
     \IT{\ELSE}{\BEGIN
                   \CASES
                    \IT{m^2 > n}{\CALL{shrink2}{n, l, m}}
                    \IT{\ELSE}{\CALL{shrink2} {n, m, u}}
                    \ENDCASES
                   \END}
                 \WHERE m = (l+u)/2
      \ENDCASES
    \END\\
    \COMMENT{ensures}{$0\leq l^2\leq n<(l+1)^2$}
  \ENDFUN
   \IN
  \BEGIN
    \CASES
    \IT{n < 0}{\bot}
    \IT{n \leq 1}{n}
    \IT{\ELSE}{\CALL{shrink2}{n,0,n}}
    \ENDCASES
  \END
  \END
 \end{pseudocode}  
\end{verbatim}

\subsection{Mixing Functional \& Imperative Styles}

Sometimes it is imperative in the interest of efficiency, to use
imperative constructs esp. when memoisation is required so that
computations are not duplicated.  Here is a simple  algorithm
for recursively computing fibonacci numbers by computing each
fibonacci number only once. Here again we have borrowed from SML
syntax.
\begin{itemize}
\item $\#1$ stands for the first component of a tuple 
\item $::$ is the {\em cons} operation on lists
\item $hd (L)$ is the first element of a non-null list
\item $tl(L)$ is the list obtained by deleting the first element
\end{itemize}

\begin{pseudocode}{fibonacci}{n}
  \BEGIN
    \GLOBAL knownFibs = [(1,1), (0,0)]\\
    \LET
    KF \GETS knownFibs\\
    \REC \FUN
       {isknown} {n, KF}
       \BEGIN
          \CASES
          \IT{KF = []}{\bot}
          \IT{hd(KF) = (n, fn)}{fn}
          \IT{\#1 (hd(KF))\neq n}{\CALL{isknown}{n, tl(KF)}}
          \ENDCASES
       \END
    \ENDFUN
    \REC \FUN
         {fib} {n}
         \BEGIN
           \CASES
           \IT {\CALL{isknown}{n, knownFibs} = (n, fn)}{fn}
           \IT {\CALL{isknown}{n, knownfibs} = \bot}
               {\BEGIN
                fn \GETS \CALL{fib}{n-1} + \CALL{fib}{n-2}; \\ 
                knownFibs := (n,fn)::knownFibs; \\
                fn
                \END}
           \ENDCASES
         \END
         \ENDFUN
   \IN\\
   \CALL{fib}{n}
  \END
\end{pseudocode}

Notice that the last clause of the definition of $\textsc{fib}$
consists of 3 lines of which the first two lines are imperative
statements which update state and the last line is the value of the
function. The above algorithm was generated by the following
\LaTeX\ code.

\begin{verbatim}
\begin{pseudocode}{fibonacci}{n}
    \GLOBAL knownFibs = [(0,0), (1,1)]\\

    \LET
    KF \GETS knownFibs\\
    \REC \FUN
       {isknown} {n, KF}
       \BEGIN
          \CASES
          \IT{KF = []}{\bot}
          \IT{hd(KF) = (n, fn)}{fn}
          \IT{\#1 (hd(KF))\neq n}{\CALL{isknown}{n, tl(KF)}}
          \ENDCASES
       \END
    \ENDFUN
    \REC \FUN
         {fib} {n}
         \BEGIN
           \CASES
           \IT {\CALL{isknown}{n, knownFibs} = (n, fn)}{fn}
           \IT {\CALL{isknown}{n, knownfibs} = \bot}
               {\BEGIN
                fn \GETS \CALL{fib}{n-1} + \CALL{fib}{n-2}; \\ 
                knownFibs := (n,fn)::knownFibs; \\
                fn
                \END}
           \ENDCASES
         \END
         \ENDFUN
   \IN\\
   \CALL{fib}{n}
  \END
\end{pseudocode}
  
\end{verbatim}

\subsection{Other Predefined Keywords}

Several other predefined keywords are available.
We summarize their usage in Table \ref{otherPredefinedKeywords}.

\begin{table}[h]
\caption{Other Predefined Keywords}
\label{otherPredefinedKeywords}
\begin{center}
\begin{tabular}{@{}ll@{}}
\hline
command & output\\
\hline
\verb+\LOCAL{list of variables}+&
$\LOCAL{\mbox{list of variables}}$\\
\verb+\GLOBAL{list of variables}+&
$\GLOBAL{\mbox{list of variables}}$\\
\verb+\EXTERNAL{list of procedures}+&
$\EXTERNAL{\mbox{list of procedures}}$\\
\verb+\RETURN{list of values}+&
$\RETURN{\mbox{list of values}}$\\
\verb+\OUTPUT{list of values}+&
$\OUTPUT{\mbox{list of values}}$\\
\verb+\EXIT+&
$\EXIT$\\
\verb+\AND+&
$\AND$\\
\verb+\OR+&
$\OR$\\
\verb+\NOT+&
$\NOT$\\
\verb+\TRUE+&
$\TRUE$\\
\verb+\FALSE+&
$\FALSE$\\
\verb+\GETS+&
$\GETS$\\
\hline
\end{tabular}
\end{center}
\end{table}

Also note 
that all of the keywords 
\verb+\IF+, \verb+\WHILE+, \verb+\CALL{}{}+, \verb+\NOT+, etc.\
are available for
use outside of the \keyword{pseudocode} environment, but they must
be input in math mode. For example,
\begin{center}
\verb+The  $\WHILE$ loop is our friend.+
\end{center}
generates
\begin{center}
The  $\WHILE$ loop is our friend.
\end{center}


\subsection{Statement Numbering}

Statements can be numbered and given a reference key so that
they can be referenced in a \LaTeX\ document using a \verb+\ref{}+ command
(see section 4.2 of \cite{lamport}).
This is done as follows:

\begin{verbatim}
\STMTNUM{<space>}{<key>}
\end{verbatim}

\noindent{}The argument \verb+<space>+ is the amount of space to be left between
the text and the statement number.  This is a length that is specified by the
user, and generally will require some experimentation in order for it to look nice.
The argument \verb+<key>+ is the reference key
used in the LaTeX\  \verb+\ref{}+ command to refer to the given
statement.

The default numbering for statements is arabic. However, it can be changed
by a suitable \verb+\renewcommand{}+.
An example is provided in the next section.

\section{An example}

The following example demonstrates the use of the
\keyword{pseudocode} environment to describe a complete
algorithm, the familiar ``mergesort'' algorithm, first in the
functional style as a higher order function on lists of elements,
where $R$ is a comparison relation on the elements of the list.
We have borrowed the list notation from ML-like languages.

The \LaTeX\ input

\begin{verbatim}
\begin{pseudocode}{mergeSort}{R, L}
\label{mergeSort}
\BEGIN
  \COMMENT{comment}{Sort a list $L$ by the ordering relation $R$}\\
  \LET 
   \REC \FUN
   {split} {L}
   \BEGIN
     \CASES
     \IT {L = []} {([], [])} 
     \IT {L = [h]} {([h], [])} 
     \IT {L = h_1::h_2::M} {(h_1::left, h_2::right)} 
     \ENDCASES\\
     \WHERE\\
     (left, right) = \CALL {split} {M}
   \END
   \ENDFUN\\
  \REC \FUN 
  {merge} {R, L, M}
    \BEGIN
      \CASES
      \IT {L = [] = M} {[]} 
      \IT {L = [] \neq M} {M}  
      \IT {L \neq [] = M} {L}  
      \IT {L \neq [] \neq M}{\BEGIN
                             \CASES
                             \IT {R(hd(L), hd(M))}
                                 {hd(L)::\CALL{merge}{tl(L), M}} 
                             \IT {\ELSE} {hd(M)::\CALL{merge}{L, tl(M)}}
                             \ENDCASES
                             \END} 
      \ENDCASES
    \END
    \ENDFUN\\
\IN 
  \BEGIN
    \CASES
    \IT {L = []} {[]} 
    \IT {L = [h]} {[h]} 
    \IT {|L|>1} {\CALL{merge} {R, sortedLeft, sortedRight}}
    \ENDCASES
  \END\\
  \WHERE\\
  (leftHalf, rightHalf) = \CALL {split} {L} \STMTNUM{1cm}{st.0}\\
  sortedLeft = \CALL {mergeSort} {R, leftHalf}\\
  sortedRight = \CALL {mergeSort} {R, rightHalf}\\
\END
\end{pseudocode}
\end{verbatim}

\noindent{}produces the following output

\begin{pseudocode}{mergeSort}{R, L}
\label{mergeSort}
\BEGIN
  \COMMENT{comment}{Sort a list $L$ by the ordering relation $R$}\\
  \LET 
   \REC \FUN
   {split} {L}
   \BEGIN
     \CASES
     \IT {L = []} {([], [])} 
     \IT {L = [h]} {([h], [])} 
     \IT {L = h_1::h_2::M} {(h_1::left, h_2::right)} 
     \ENDCASES\\
     \WHERE\\
     (left, right) = \CALL {split} {M}
   \END
   \ENDFUN\\
  \REC \FUN 
  {merge} {R, L, M}
    \BEGIN
      \CASES
      \IT {L = [] = M} {[]} 
      \IT {L = [] \neq M} {M}  
      \IT {L \neq [] = M} {L}  
      \IT {L \neq [] \neq M}{\BEGIN
                             \CASES
                             \IT {R(hd(L), hd(M))}
                                 {hd(L)::\CALL{merge}{tl(L), M}} 
                             \IT {\ELSE} {hd(M)::\CALL{merge}{L, tl(M)}}
                             \ENDCASES
                             \END} 
      \ENDCASES
    \END
    \ENDFUN\\
\IN 
  \BEGIN
    \CASES
    \IT {L = []} {[]} 
    \IT {L = [h]} {[h]} 
    \IT {|L|>1} {\CALL{merge} {R, sortedLeft, sortedRight}}
    \ENDCASES
  \END\\
  \WHERE\\
  (leftHalf, rightHalf) = \CALL {split} {L} \STMTNUM{1cm}{st.0}\\
  sortedLeft = \CALL {mergeSort} {R, leftHalf}\\
  sortedRight = \CALL {mergeSort} {R, rightHalf}\\
\END
\end{pseudocode}


In particular notice the use of the equality symbol in (\ref{st.0})
and the statements following it, as distinct from the $\DEF$ used in the
definition of the functions. Further also note the multiple occurrences of the variable name $L$ in the different scopes. The scopes are clarified by the use of braces in nested forms. 

\noindent{}The following \LaTeX\ source describes the mergesort
algorithm in the imperative style on arrays of integers
under the usual $<$ ordering. The \LaTeX\ input

\begin{verbatim}
\renewcommand{\thepseudonum}{\roman{pseudonum}}
\begin{pseudocode}{MergeSort}{n,X}
\label{MergeSort}
\COMMENT{comment}{Sort the array $X$ of length $n$}\\
\IF n=2 \THEN
\BEGIN
  \IF X[0]>X[1] \THEN
  \BEGIN
    T \GETS X[0]\\
    X[0]\GETS X[1]\\
    X[1]\GETS T
  \END
\END
\ELSEIF n>2 \THEN
\BEGIN
  m\GETS \lfloor n/2 \rfloor\\
  \FOR i\GETS 0 \TO m-1 \DO A[i] \GETS X[i]\\
  \FOR i\GETS m \TO n-1 \DO B[i] \GETS X[i]\\
  \COMMENT{comment}{Now sort the subarrays $A$ and $B$}\\
  \CALL{MergeSort}{m,A}\\
  \CALL{MergeSort}{n-m,B}\\
  i\GETS 0\\
  j\GETS 0\\
  \FOR k \GETS 0 \TO n-1 \DO
  \BEGIN
    \IF A[i] \leq B[j] \THEN
    \BEGIN
       X[k]\GETS A[i] \STMTNUM{1in}{st.1}\\
       i\GETS i+1
    \END
    \ELSE
    \BEGIN
       X[k]\GETS B[j] \STMTNUM{1.03in}{st.2}\\
       j\GETS j+1
    \END
  \END
\END
\end{pseudocode}
\end{verbatim}

\noindent{}produces the following output:

\renewcommand{\thepseudonum}{\roman{pseudonum}}
\begin{pseudocode}{MergeSort}{n,X}
\label{MergeSort}
\COMMENT{comment}{Sort the array $X$ of length $n$}\\
\IF n=2 \THEN
\BEGIN
  \IF X[0]>X[1] \THEN
  \BEGIN
    T \GETS X[0]\\
    X[0]\GETS X[1]\\
    X[1]\GETS T
  \END
\END
\ELSEIF n>2 \THEN
\BEGIN
  m\GETS \lfloor n/2 \rfloor\\
  \FOR i\GETS 0   \TO m-1 \DO A[i] \GETS X[i]\\
  \FOR i\GETS m \TO n-1 \DO B[i] \GETS X[i]\\
  \COMMENT{comment}{Now sort the subarrays $A$ and $B$}\\
  \CALL{MergeSort}{m,A}\\
  \CALL{MergeSort}{n-m,B}\\
  i\GETS 0\\
  j\GETS 0\\
  \FOR k \GETS 0 \TO n-1 \DO
  \BEGIN
    \IF A[i] \leq B[j] \THEN
    \BEGIN
       X[k]\GETS A[i] \STMTNUM{.9in}{st.1}\\
       i\GETS i+1
    \END
    \ELSE
    \BEGIN
       X[k]\GETS B[j] \STMTNUM{.93in}{st.2}\\
       j\GETS j+1
    \END
  \END
\END
\end{pseudocode}

The counter \verb+pseudonum+ keeps track of the statement numbers.
The style of the counter values can be changed using the method
described in Section 6.3 of \cite{lamport}. For example,
we used the command
\begin{center}
\verb+ \renewcommand{\thepseudonum}{\roman{pseudonum}} +
\end{center}
in our example so that statements were numbered with 
lowercase Roman numerals.
We also assigned a label to the algorithm using the
\verb+\label{}+ command that is described in Section 4.2 of \cite{lamport}.
Finally, by trial and error, we determined 
spacing so that the statement numbers would be vertically aligned.

We now give an example of how the numbered statements in the above
algorithm can be referenced in a \LaTeX document.  The commands

\begin{verbatim}
On lines (\ref{st.1}) and (\ref{st.2}) of Algorithm 
\ref{MergeSort}, we determine the $k$th element of the 
sorted array.
\end{verbatim}

\noindent{}produce the following output:

\medskip 

\noindent{}On lines (\ref{st.1}) and (\ref{st.2}) of Algorithm \ref{MergeSort},  
we determine the $k$th element of the 
sorted array.

%\medskip 



\section{Framing}
The \keyword{pseudocode} environment also has an optional parameter, 
\verb+<frame>+. The complete form of the \keyword{pseudocode} environment is

\begin{verbatim}
\begin{pseudocode}[<frame>]{<Name>}{<Parameters>}
  pseudocode constructs
\end{pseudocode}
\end{verbatim}

The possible values of  \verb+<frame>+ are:
\begin{center}
\begin{tabular}{llll}
\verb+shadowbox+ &
\verb+doublebox+ &
\verb+ovalbox+ &
\verb+Ovalbox+  \\
\verb+framebox+ &
\verb+plain+ &
\verb+ruled+ &
\verb+display+
\end{tabular}
\end{center}
The values ending with ``box'' draw various types of frames around the algorithm.
The value \verb+plain+ is the default and adds no frame to the algorithm.
The value \verb+display+ is used for displaying sections of code
without the algorithm name or parameters. Here are some examples with
input:

\begin{verbatim}
\begin{pseudocode}[<frame>]{SquareAndMultiply}{x,b,n}
 \COMMENT{comment}{ Compute $x^b \pmod{n}$}\\
 z\GETS 1\\
 \WHILE b > 0 \DO
 \BEGIN
  z \GETS z^2 \pmod{n} \\
  \IF b\mbox{ is odd} 
    \THEN z \GETS z \cdot x  \pmod{n} \\
    b \GETS \CALL{ShiftRight}{b}
 \END\\
 \RETURN{z}
\end{pseudocode}
\end{verbatim}

\noindent{}where we give \verb+<frame>+  each of the values described above.

\medskip

\noindent When \verb+<frame>+ is \verb+shadowbox+ we obtain:

\begin{pseudocode}[shadowbox]{SquareAndMultiply}{x,b,n}
 \COMMENT{comment}{ Compute $x^b \pmod{n}$}\\
 z\GETS 1\\
 \WHILE b > 0 \DO
 \BEGIN
  z \GETS z^2 \pmod{n} \\
  \IF b\mbox{ is odd} 
    \THEN z \GETS z \cdot x  \pmod{n} \\
    b \GETS \CALL{ShiftRight}{b}
 \END\\
 \RETURN{z}
\end{pseudocode}

\noindent When \verb+<frame>+ is \verb+doublebox+ we obtain:

\begin{pseudocode}[doublebox]{SquareAndMultiply}{x,b,n}
 \COMMENT{comment}{ Compute $x^b \pmod{n}$}\\
 z\GETS 1\\
 \WHILE b > 0 \DO
 \BEGIN
  z \GETS z^2 \pmod{n} \\
  \IF b\mbox{ is odd} 
    \THEN z \GETS z \cdot x  \pmod{n} \\
    b \GETS \CALL{ShiftRight}{b}
 \END\\
 \RETURN{z}
\end{pseudocode}

\noindent When \verb+<frame>+ is \verb+ovalbox+ we obtain:

\begin{pseudocode}[ovalbox]{SquareAndMultiply}{x,b,n}
 \COMMENT{comment}{ Compute $x^b \pmod{n}$}\\
 z\GETS 1\\
 \WHILE b > 0 \DO
 \BEGIN
  z \GETS z^2 \pmod{n} \\
  \IF b\mbox{ is odd} 
    \THEN z \GETS z \cdot x  \pmod{n} \\
    b \GETS \CALL{ShiftRight}{b}
 \END\\
 \RETURN{z}
\end{pseudocode}

\noindent When \verb+<frame>+ is \verb+Ovalbox+ we obtain:

\begin{pseudocode}[Ovalbox]{SquareAndMultiply}{x,b,n}
 \COMMENT{comment}{ Compute $x^b \pmod{n}$}\\
 z\GETS 1\\
 \WHILE b > 0 \DO
 \BEGIN
  z \GETS z^2 \pmod{n} \\
  \IF b\mbox{ is odd} 
    \THEN z \GETS z \cdot x  \pmod{n} \\
    b \GETS \CALL{ShiftRight}{b}
 \END\\
 \RETURN{z}
\end{pseudocode}

\noindent When \verb+<frame>+ is \verb+framebox+ we obtain:

\begin{pseudocode}[framebox]{SquareAndMultiply}{x,b,n}
 \COMMENT{comment}{ Compute $x^b \pmod{n}$}\\
 z\GETS 1\\
 \WHILE b > 0 \DO
 \BEGIN
  z \GETS z^2 \pmod{n} \\
  \IF b\mbox{ is odd} 
    \THEN z \GETS z \cdot x  \pmod{n} \\
    b \GETS \CALL{ShiftRight}{b}
 \END\\
 \RETURN{z}
\end{pseudocode}

\noindent When \verb+<frame>+ is \verb+plain+ or if \verb+[<frame>]+ is omitted we obtain:

\begin{pseudocode}[plain]{SquareAndMultiply}{x,b,n}
 \COMMENT{comment}{ Compute $x^b \pmod{n}$}\\
 z\GETS 1\\
 \WHILE b > 0 \DO
 \BEGIN
  z \GETS z^2 \pmod{n} \\
  \IF b\mbox{ is odd} 
    \THEN z \GETS z \cdot x  \pmod{n} \\
    b \GETS \CALL{ShiftRight}{b}
 \END\\
 \RETURN{z}
\end{pseudocode}

\noindent When \verb+<frame>+ is \verb+ruled+ we obtain:

\begin{pseudocode}[ruled]{SquareAndMultiply}{x,b,n}
 \COMMENT{comment}{ Compute $x^b \pmod{n}$}\\
 z\GETS 1\\
 \WHILE b > 0 \DO
 \BEGIN
  z \GETS z^2 \pmod{n} \\
  \IF b\mbox{ is odd} 
    \THEN z \GETS z \cdot x  \pmod{n} \\
    b \GETS \CALL{ShiftRight}{b}
 \END\\
 \RETURN{z}
\end{pseudocode}


The purpose of the value \verb+display+ is to allow portions of
algorithms to be displayed with out the algorithm header. Thus
for example to display the section of code in the $\WHILE$ loop
of the $\CALL{SquareAndMultiply}{}$ algorithm one could write

\begin{verbatim}
\begin{center}
\begin{minipage}{2in}
\begin{pseudocode}[display]{}{}
  z \GETS z^2 \pmod{n} \\
  \IF b\mbox{ is odd}
    \THEN z \GETS z \cdot x  \pmod{n}\\
    b \GETS \CALL{ShiftRight}{b}
\end{pseudocode}
\end{minipage}
\end{center}
\end{verbatim}

\noindent which would produce the following output:

\begin{center}
\begin{minipage}{2in}
\begin{pseudocode}[display]{}{}
  z \GETS z^2 \pmod{n} \\
  \IF b\mbox{ is odd}
    \THEN z \GETS z \cdot x  \pmod{n}\\
     b \GETS \CALL{ShiftRight}{b}
\end{pseudocode}
\end{minipage}
\end{center}

\section{Logistics \& pragmatics}

The file \keyword{funpseudocode.sty} may be freely downloaded from 

\centerline{\url{http://www.cse.iitd.ac.in/~sak/funpseudocode.sty}}

\noindent{}and copied into the \keyword{pseudocode} folder (alongside the older
\keyword{pseudocode.sty} file) so that \keyword{pdflatex} can find it. In my
Ubuntu system, I have copied it into the folder

\centerline{\keyword{/usr/share/texlive/texmf-dist/tex/latex/pseudocode/}}
\noindent{}where it sits alongside the older \keyword{pseudocode.sty} file.
Then  you may just include 

\centerline{{\tt \textbackslash usepackage\{funpseudocode\}}} 

\noindent{}in the preamble of your \LaTeX\ file to invoke its commands in the
environment {\tt pseudocode}. Do not put both {\tt \textbackslash
  usepackage\{funpseudocode\}} and {\tt \textbackslash usepackage\{pseudocode\}}
in your preamble.
\begin{thebibliography}{99}

\bibitem{LC}
M.\ Goossens, F.\ Mittelbach and A.\ Samarin, {\it The \LaTeX Companion}, 
Addison-Wesley, 1994.

\bibitem{cages}
D.L.\ Kreher and D.R.\ Stinson,
{\em Combinatorial Algorithms: Generation, Enumeration and Search},
CRC Press, 1999.

\bibitem{lamport}
L.\ Lamport, \LaTeX, 
{\em A Document Preparation System},
Addison-Wesley, 1994.

\bibitem{lehrer}
T.\ Lehrer,
{\tt
  \url{http://www.lyricsfreak.com/t/tom+lehrer/lobachevsky\_20138409.html}}
20th century.
\bibitem{lloyd}
J.\ W.\ Lloyd,
{\em Foundations of Logic Programming},
Springer-Verlag, 1987.
\bibitem{backus}
J.\ Backus,
{\it Can Programming Be Liberated from the von Neumann Style? A
  Functional Style and Its Algerbra of Programs},
CACM, 1978.

\end{thebibliography}
\end{document}

